\begin{saresult}{A.128}{Known}{The three- and four-sphere component cases follow directly in PL}
For every standard PL sphere $S^n$, $n\geq2$,
\[
 \pi_0 P_{\mathrm{PL}}(S^n)=0.
\]
Consequently the same assertion holds for each definable model of
the standard PL sphere over every real closed field, with the
simplicial concordance convention of this study. In particular,
this closes the standard $S^3$ and $S^4$ component cases without
a smooth-to-PL comparison in base dimension three or four.
\end{saresult}
\begin{proof}
Write
\[
 X_j=\mathrm{PL}_j/\mathrm O_j,
\]
using the simplicial groups of germs and their homogeneous spaces.
These are the unstable smoothing spaces, not their stable
substitutes. The sphere calculation in
\cite[\S6, p.~43, equation~(1), following Proposition~5.3]{BurgheleaLashof1974II}
gives a weak equivalence of spaces
\[
 P_{\mathrm{PL}}(S^n)
 \simeq E(X_{n+1},X_n)
 :=\operatorname{hofib}_{*}(X_n\longrightarrow X_{n+1}).
                                                        \tag{A.128.1}
\]
Here $E(X_{n+1},X_n)$ is the path space defined in the same
paper, p.~42: a point consists of a point of $X_n$ and a path
in $X_{n+1}$ from the basepoint to its image. The sphere calculation
uses the Kuiper--Lashof fibration, which the paper states on
p.~38 before Proposition~5.3 as
\[
 \mathrm{PL}_n\longrightarrow\mathrm{PL}_{n+1}
 \longrightarrow S^n\times B P_{\mathrm{PL}}(S^n),
\]
with the orthogonal sphere fibration mapping into it. It has no
high-dimensional stability hypothesis. The original
\cite[Theorem~1 and definitions, pp.~243--244]{KuiperLashof1966II}
defines each simplex by a PL bundle isomorphism on the entire
product $\Delta_k\times X$, with base $\Delta_k$; its concordance
group fixes $S^n\times0$ pointwise and leaves the top free.
On these compact products the PL maps have finite triangulations,
so the convention is precisely the jointly PL convention used here.
Only path components of
(A.128.1) will be used, so no compatibility of this equivalence
with a proposed multiplication is required.

The low-degree stabilization statement
\cite[\S5, p.~37, formula~(3)]{BurgheleaLashof1974II}
is
\[
 \pi_i(X_j)\xrightarrow{\ \cong\ }
 \pi_i(\mathrm{PL}/\mathrm O)\qquad(i<j).
                                                        \tag{A.128.2}
\]
The stable space $\mathrm{PL}/\mathrm O$ is connected and has
vanishing homotopy groups through degree six; the relevant
vanishing is explicitly recorded in
\cite[Annex~C, p.~318, calculation following Theorem~5.4]{KirbySiebenmann1977}.
For $n\geq2$, (A.128.2) therefore gives
\[
 \pi_0(X_n)=0,\qquad \pi_1(X_{n+1})=0.
\]
The exact homotopy sequence of (A.128.1) now contains
\[
 \pi_1(X_{n+1})\longrightarrow
 \pi_0 P_{\mathrm{PL}}(S^n)\longrightarrow\pi_0(X_n),
\]
and both outer terms are trivial. The middle pointed set has
one element, hence the concordance component group is trivial.

For a general real closed field, A.115 transports the real PL
component computation. The relative one-parameter component
comparison in A.113 transports it to the stipulated definable
group. A semialgebraic sphere may be replaced by its finite PL
model before this comparison.

There are two distinct statements in the source that must not be
confused here. Formula~(3) on p.~37 is the low-degree assertion
$i<j$ used above. The stronger Theorem~5.1 on that page is a
high-dimensional stability assertion and is not used. Likewise,
Theorem~6.1 on p.~41, whose printed statement omits a dimension
restriction although its displayed proof invokes Theorem~5.1,
is not used in base dimension three or four. The argument gives
no component computation for a general closed three- or
four-manifold.
\end{proof}
